% =========================================================
\subsection{Resident / Public}
% =========================================================
For residents and the public, the questions are \emph{what the model pays attention to across
Hull}, and \emph{what it weighs in a place like theirs} (Figure~\ref{fig:stakeholder-public}).\par\medskip
% Resident / Public stakeholder card (Chapter 6). Same design as the other cards: two dashboard
% panels, four key figures, two text columns. Example area: Stoneferry 016B (E01012806), a
% typical area (81st of 168, total importance close to the Hull median), not the rank-1 extreme.
% Plain language throughout: the panels use the dashboard's plain public view (no coefficients
% or model terms), and each figure appears once (the panels print the area's rank and its top
% shares, so tiles and text do not repeat them; checked by card_mode).
% Every image is the dashboard's own short-card public panel, drawn by
% dashboard.views.stakeholder_figures and exported at print size by
%   python export_dashboard_panels.py --stakeholder public_short --lsoa E01012806
%   01 where predictions depend most on an area's own data   02 what the model weighs in Stoneferry.
% Text and tiles: generated by the dashboard LLM pipeline (generate_card_narrative.py
% --stakeholder public --openai-model gpt-4.1): master prompt geosmile-hull-v8 + public card-mode
% addendum. The model marks figures and key terms itself. Prompts, evidence, raw output and
% checks: public_card_llm.json. The generated headline is not shown on the card.
%
% Preamble: \usepackage[most]{tcolorbox}, xcolor, graphicx.
% Copy into Chapter6/figs/: dashboard_public_short_E01012806_{1,2}.png and
% public_card_llm_content_openai.tex.
% =========================================================

\definecolor{cardInk}{HTML}{12171C}
\definecolor{cardMuted}{HTML}{57646D}
\definecolor{cardTeal}{HTML}{1F4E5F}
\definecolor{cardPink}{HTML}{FCE4EC}
\definecolor{cardWine}{HTML}{8A0033}
\definecolor{cardRule}{HTML}{D9DEE2}

% Same macros as planner_card.tex (\providecommand: defined once, whichever card comes first).
\providecommand{\hlpink}[1]{\begingroup\setlength{\fboxsep}{0.6pt}\colorbox{cardPink}{\textcolor{cardWine}{#1}}\endgroup}
\providecommand{\pv}[1]{\hlpink{#1}}
\providecommand{\kt}[1]{\hlpink{#1}}
\providecommand{\roletag}[1]{\tcbox[on line, colback=cardTeal, colframe=cardTeal, arc=2.5mm,
  boxsep=0pt, left=6pt, right=6pt, top=3pt, bottom=3pt]{\color{white}\sffamily\bfseries #1}}
\providecommand{\kicker}[1]{{\color{cardMuted}\sffamily\small\bfseries #1}}
\providecommand{\tile}[2]{\begin{tcolorbox}[enhanced, nobeforeafter, width=\cardtilewidth,
  equal height group=\cardtilegroup, colback=cardPink, frame hidden, arc=2mm,
  left=5pt, right=3pt, top=3pt, bottom=3pt, valign=top, before upper=\raggedright]
  {\bfseries\color{cardWine}#1}\par\vspace{1pt}
  {\sffamily\color{cardMuted}#2}\end{tcolorbox}}
\providecommand{\dashpanel}[3]{\kicker{#1}\par\vspace{2pt}%
  {\centering\includegraphics[width=#2\linewidth]{#3}\par}}

\def\cardtilegroup{publictiles}
\def\cardtilewidth{0.2375\linewidth}
\providecommand{\cardHeadline}{[headline]}
\providecommand{\cardSecA}{[section 01]}
\providecommand{\cardSecB}{[section 02]}
\providecommand{\cardFooter}{[footer]}
\providecommand{\cardModel}{[model]}
\providecommand{\cardTileA}{\tile{--}{--}}
\providecommand{\cardTileB}{\tile{--}{--}}
\providecommand{\cardTileC}{\tile{--}{--}}
\providecommand{\cardTileD}{\tile{--}{--}}
\input{Chapter6/figs/public_card_llm_content_openai.tex}

\begin{figure}[!htbp]
\begin{tcolorbox}[enhanced, colback=white, colframe=cardRule, arc=3mm, boxrule=0.6pt,
                  left=3mm, right=3mm, top=2.5mm, bottom=2.5mm, before skip=0pt, after skip=2pt,
                  fontupper=\sffamily\normalsize\color{cardInk}]
\setlength{\parskip}{0pt}\raggedright%
\roletag{RESIDENT / PUBLIC}\par\smallskip
{\color{cardMuted}What the model pays attention to, across Hull and in a typical area}\par\vspace{4pt}
% Row 1: the two panels side by side.
\begin{minipage}[t]{0.45\linewidth}\vspace{0pt}
  \dashpanel{01 \textperiodcentered{} Across Hull}{1.0}{Chapter6/figs/dashboard_public_short_E01012806_1.png}
\end{minipage}\hfill
\begin{minipage}[t]{0.52\linewidth}\vspace{0pt}
  \dashpanel{02 \textperiodcentered{} A typical area: Stoneferry}{1.0}{Chapter6/figs/dashboard_public_short_E01012806_2.png}
\end{minipage}\par\vspace{4pt}
% Row 2: key figures.
\kicker{Key figures}\par\vspace{2pt}
\noindent\cardTileA\hfill\cardTileB\hfill\cardTileC\hfill\cardTileD\par\vspace{6pt}
% Row 3: the two text sections side by side.
\begin{minipage}[t]{0.485\linewidth}\vspace{0pt}\raggedright
  \kicker{Across Hull}\par\vspace{2pt}
  \cardSecA\par
\end{minipage}\hfill
\begin{minipage}[t]{0.485\linewidth}\vspace{0pt}\raggedright
  \kicker{In Stoneferry}\par\vspace{2pt}
  \cardSecB\par
\end{minipage}\par\vspace{4pt}
{\color{cardRule}\rule{\linewidth}{0.4pt}}\par\vspace{1pt}
{\color{cardMuted}\footnotesize \cardFooter{} Panels are the dashboard's public view. Text, tiles
and \pv{highlights} generated by \cardModel{} (prompt geosmile-hull-v8, card mode).\par}
\end{tcolorbox}
\caption{Resident / Public card: where predictions depend most on an area's own data, what the
model weighs in a typical area (Stoneferry), and the LLM-generated plain-language explanation.}
\label{fig:stakeholder-public}
\end{figure}
